\documentclass[14pt, a4paper]{extreport}

\usepackage{kpi-lab}
% Defined by subject
\newcommand{\CourseTitle}{Проектування програмних систем для мобільних пристроїв}
\newcommand{\Teacher}{асистент, Нестерук А. О.}
\newcommand{\ReportYear}{2026}
% Defined by Student
\newcommand{\StudentName}{Ковальов Олександр}
\newcommand{\StudentGroup}{ІМ-51мн}
\newcommand{\Variant}{4}
\newcommand{\CourseNumber}{2}
% Defined by Submission
\newcommand{\Topic}{Архітектура додатків та допоміжні бібліотеки}
\newcommand{\LabNumber}{5}

\begin{document}
	\pagenumbering{gobble}
	
	\begin{titlepage}
		\begin{center}
			\normalsize
			Національний технічний університет України\\
			«Київський політехнічний інститут імені Ігоря Сікорського» \\[0.5em]
			Факультет інформатики та обчислювальної техніки\\
			Кафедра обчислювальної техніки
			
			\vfill
			
			{\large \textbf{ЗВІТ}\\[0.5em]}
			{з лабораторної роботи №\LabNumber \\}
			{з дисципліни <<\CourseTitle>>}
			
			\vspace{2em}
			
			{\textbf{<<\Topic>>}}
			
			\vspace{1em}
			
			{\textbf{Варіант №\Variant}}
			
			\vfill
			
			\begin{flushright}
				Виконав: \\ студент \CourseNumber{} курсу, групи \StudentGroup \\
				\StudentName \\[1em]
				Перевірив: \\ \Teacher
			\end{flushright}
			
			\vfill
			
			\begin{flushright}
				Дата здачі: \DTMtoday
			\end{flushright}
			
			\vspace{5em}
			КИЇВ -- \ReportYear
		\end{center}
	\end{titlepage}
	
	\clearpage
	\pagenumbering{arabic}
	\setcounter{page}{2}
	
	\subsection*{Мета} 
	
	Вивчити особливості архітектури програм для операційної системи Android та виділити їх як для платформи розробки програм.
	
	\subsection*{Завдання}
	
	Додати у програму, розроблену у лабораторній роботі №4, сторонні бібліотеки, різними методами. Зробіть висновок про зручність цих методів. Скористайтеся бібліотеками з відкритих джерел Інтернету.
	
	\subsection*{Завдання за варіантом}
	
	Доменна область: Абітурієнт. 
	
	Поля:  ID, Прізвище, Ім'я, По-батькові, Адреса, Телефон, Оцінки. 
	
	Створити масив об'єктів. 
	
	Вивести:
	
	\begin{enumerate}[label=\alph*.]
		\item список абітурієнтів, які мають незадовільні оцінки;
		
		\item список абітурієнтів, середній бал у яких вище заданого;
		
		\item вибрати задане число \texttt{N} абітурієнтів, що мають найбільш високій середній бал (вивести також повний список абітурієнтів, які мають напівпрохідний бал).
	\end{enumerate}
	
	\section*{Хід роботи.}

	Для виконання роботи використано застосунок, розроблений у четвертій лабораторній роботі, без зміни архітектури MVC чи стеку -- Rust, Tauri, React з TypeScript. Завдання вимагає додати до застосунку сторонні бібліотеки різними методами та зробити висновок про зручність цих методів; для порівняння обрано два принципово різні методи -- підключення через реєстр пакетів (pnpm) та підключення через git-залежність (не crates.io).
	
	Першим методом додано бібліотеку \texttt{fuse.js} -- звичайним способом, командою \texttt{pnpm add fuse.js}, яка резолвить пакет через реєстр npm, записує точну версію в \texttt{package.json} і фіксує її разом з хешем цілісності в lock-файлі.
	
	\image{3cm}{fuse}
	
	Бібліотеку використано на екрані <<Пошук>>: замість попереднього точного пошуку підрядка на боці Rust, екран тепер один раз завантажує повний список абітурієнтів і виконує нечіткий (typo-tolerant) пошук на клієнті за допомогою індексу \texttt{Fuse}, побудованого з полем \texttt{last\_name} та порогом нечіткості \texttt{threshold: 0.4}. Фільтрація відбувається одразу при зміні поля вводу, без окремої команди на сервер. Коректність перевірено на емуляторі: прізвище введено з однією помилкою (одна літера), і застосунок все одно знайшов відповідного абітурієнта.
	
	\image{4cm}{fuzzy}
	
	Другим методом додано крейт \texttt{qrcode} -- через git-залежність у \texttt{Cargo.toml}: замість номера версії з crates.io вказано посилання на GitHub-репозиторій та тег \texttt{v0.14.1}. На відміну від звичайних залежностей, Cargo в цьому випадку клонує репозиторій безпосередньо за вказаним тегом, а не отримує пакет з реєстру crates.io.
	
	\image{3cm}{cargo}
	
	На боці Rust додано контролер-модуль із відповідною командою, яка знаходить абітурієнта за \texttt{ID}, формує з його даних (ПІБ, адреса, телефон) текстовий рядок і кодує його у QR-код, після чого рендерить його в SVG-розмітку. На боці React додано кнопку <<QR-код>> у кожному рядку списку абітурієнтів на екрані <<Абітурієнти>>. За натисканням під таблицею з'являється панель з отриманим SVG.
	
	\image{8cm}{page}
	
	Коректність згенерованого QR-коду перевірено, зчитавши його через сайт -- закодований текст відповідав даним вибраного абітурієнта.
	
	\image{5cm}{qr-result}

	\subsection*{Висновок}

	У цій лабораторній роботі до застосунку <<Абітурієнт>> додано дві сторонні бібліотеки різними шляхами. 
	
	Порівнюючи зручність методів, підключення через реєстр пакетів виявилося значно простішим і надійнішим: одна команда знаходить пакет, фіксує точну версію й хеш цілісності в lock-файлі, а оновлення версії так само зводиться до однієї команди; при цьому пакет пройшов публікацію в реєстрі, що дає певну гарантію цілісності та доступності. Підключення через git-залежність вимагає більше ручної роботи й уважності: Cargo не має семантичного резолвера версій для таких залежностей, тому без явного зазначення тегу чи коміту збірка була б нестабільною -- залежність могла б непомітно змінитися разом зі змінами в репозиторії. З іншого боку, цей метод єдиний застосовний, коли бібліотека не опублікована в crates.io взагалі, або коли потрібна саме версія з конкретної гілки чи форку. Відповідно, реєстр пакетів варто використовувати як основний спосіб, а git-залежність -- лише тоді, коли реєстр не надає потрібного варіанта бібліотеки.
	
	\section*{Контрольні запитання.}

	\begin{enumerate}
		\item \textbf{Охарактеризуйте архітектуру Android-програми.} Архітектура Android-програми базується на чотирьох фундаментальних компонентах: активностях (Activities) для візуальної взаємодії з користувачем, фонових службах (Services) для виконання тривалих операцій без графічного інтерфейсу, приймачах широкомовних повідомлень (Broadcast Receivers) для реакції на загальносистемні події та контент-провайдерах (Content Providers) для управління доступом та безпечного обміну даними між різними додатками. Сучасний підхід до розробки від Google також вимагає дотримання патерну MVVM (Model-View-ViewModel), де архітектура чітко розділяється на рівень інтерфейсу (View), рівень управління станом (ViewModel) та рівень даних із репозиторіями, що забезпечує стабільність роботи при зміні конфігурації екрану та спрощує тестування коду.
		
		\item \textbf{Назвіть класи-помічники.} Класи-помічники (Helper classes або Utility classes) у розробці під Android -- це спеціальні класи, які інкапсулюють рутинну, повторювану логіку або специфічні алгоритми для розвантаження основних компонентів програми. Класичним системним прикладом є \texttt{SQLiteOpenHelper}, який бере на себе всю роботу зі створення, відкриття та управління версіями локальної бази даних, але до цієї категорії також належать будь-які власні класи зі статичними методами для валідації введених даних, форматування дат, конвертації величин, роботи з мережею або обробки зображень, які не прив'язані до життєвого циклу активностей чи фрагментів.
		
		\item \textbf{Назвіть основні етапи підключення сторонніх бібліотек у проект.} Процес інтеграції розпочинається зі знаходження необхідної залежності в публічних репозиторіях на кшталт Maven Central, після чого розробник відкриває конфігураційний файл \texttt{build.gradle} (на рівні модуля додатка) і додає рядок із вказівкою групи, назви та версії інструменту у спеціальний блок dependencies. Завершальним і найважливішим етапом є натискання кнопки <<Sync Now>> в Android Studio, що запускає процес синхронізації Gradle, під час якого система автоматично завантажує всі потрібні файли, перевіряє сумісність версій та робить нові класи доступними для використання в основному коді проекту.
	\end{enumerate}

	\pagebreak

	\section*{Лістинг.}

	\textbf{qr.rs}

	\begin{minted}{rust}
		//! Lab 5: Generating a QR code containing applicant data using the `qrcode` crate,
		//! integrated via a Git dependency (rather than crates.io) — a demonstration of an alternative
		//! method for adding third-party code to a project.
		
		// State<T> must be taken by value in every command here — that's Tauri's
		// calling convention, not a real ownership need.
		#![allow(clippy::needless_pass_by_value, reason = "Tauri's convention")]
		
		use crate::repository::Repository;
		use qrcode::QrCode;
		use qrcode::render::svg;
		use std::sync::{Mutex, PoisonError};
		use tauri::State;
		
		#[tauri::command]
		pub fn generate_qr(
		state: State<'_, Mutex<Repository>>, id: u32,
		) -> Result<String, String> {
			let (last_name, first_name, patronymic, address, phone) = state
			.lock()
			.unwrap_or_else(PoisonError::into_inner)
			.list()
			.iter()
			.find(|a| a.id == id)
			.map(|a| {
				(
				a.last_name.clone(),
				a.first_name.clone(),
				a.patronymic.clone(),
				a.address.clone(),
				a.phone.clone(),
				)
			})
			.ok_or_else(|| format!("Абітурієнта з ID {id} не знайдено"))?;
			
			let payload = format!(
			"ID: {id}\nПІБ: {last_name} {first_name} {patronymic}\nАдреса: {address}\nТелефон: {phone}",
			);
			
			let code = QrCode::new(payload.as_bytes()).map_err(|err| err.to_string())?;
			
			Ok(code.render::<svg::Color>().build())
		}
	\end{minted}
	
	\textbf{Search.tsx}
	
	\begin{minted}{tsx}
		import { useEffect, useMemo, useState } from "react";
		import { invoke } from "@tauri-apps/api/core";
		import { Link } from "react-router-dom";
		import Fuse from "fuse.js";
		import type { Applicant } from "../types";
		
		// Demonstrates Button, TextView, EditText, and ListView on one screen
		// (Lab 4 assignment). Lab 5 adds fuse.js here (via the pnpm/npm registry)
		// for typo-tolerant fuzzy matching on last names — the other new library
		// this lab adds, the `csv` crate on the Students screen, is pulled in via
		// a git dependency instead, to compare the two installation methods.
		export function Search() {
			const [query, setQuery] = useState("");
			const [all, setAll] = useState<Applicant[] | null>(null);
			const [selectedId, setSelectedId] = useState<number | null>(null);
			
			useEffect(() => {
				invoke<Applicant[]>("list_applicants").then(setAll);
			}, []);
			
			const fuse = useMemo(
			() => new Fuse(all ?? [], { keys: ["last_name"], threshold: 0.4 }),
			[all],
			);
			
			const rows = useMemo(() => {
				if (!all) {
					return null;
				}
				if (query.trim() === "") {
					return all;
				}
				return fuse.search(query).map((result) => result.item);
			}, [all, query, fuse]);
			
			function clear() {
				setQuery("");
				setSelectedId(null);
			}
			
			const selected = rows?.find((a) => a.id === selectedId) ?? null;
			
			return (
			<main className="container">
			<Link className="back" to="/">
			← На головну
			</Link>
			<h1>Пошук абітурієнтів</h1>
			<p className="subtitle">
			EditText + Button + ListView + TextView, нечіткий пошук через fuse.js
			</p>
			
			{/* EditText + Button */}
			<form className="controls" onSubmit={(e) => e.preventDefault()}>
			<label>
			Прізвище (допускає помилки)
			<input
			type="text"
			value={query}
			onChange={(e) => setQuery(e.currentTarget.value)}
			/>
			</label>
			<button type="button" onClick={clear}>
			Очистити
			</button>
			</form>
			
			{/* ListView */}
			<section className="panel">
			<h2>
			Результати <span className="count">({rows?.length ?? 0})</span>
			</h2>
			{rows && rows.length > 0 ? (
				<ul className="list-view">
				{rows.map((a) => (
					<li
					key={a.id}
					className={a.id === selectedId ? "selected" : undefined}
					onClick={() => setSelectedId(a.id)}
					>
					{a.last_name} {a.first_name} {a.patronymic} — сер. бал{" "}
					{a.average}
					</li>
					))}
				</ul>
				) : (
				<p className="empty">— немає —</p>
				)}
			</section>
			
			{/* TextView */}
			<section className="panel">
			<h2>Деталі</h2>
			{selected ? (
				<p className="text-view">
				ID: {selected.id}
				<br />
				ПІБ: {selected.last_name} {selected.first_name}{" "}
				{selected.patronymic}
				<br />
				Адреса: {selected.address}
				<br />
				Телефон: {selected.phone}
				<br />
				Оцінки: {selected.grades.join(", ")}
				<br />
				Середній бал: {selected.average}
				</p>
				) : (
				<p className="empty">Виберіть абітурієнта зі списку</p>
				)}
			</section>
			</main>
			);
		}
	\end{minted}
	
	\textbf{Students.tsx}
	
	\begin{minted}{tsx}
		import { useEffect, useState, type FormEvent } from "react";
		import { invoke } from "@tauri-apps/api/core";
		import { Link } from "react-router-dom";
		import type { Applicant, ApplicantInput } from "../types";
		
		const emptyForm = {
			last_name: "",
			first_name: "",
			patronymic: "",
			address: "",
			phone: "",
			gradesText: "",
		};
		
		function parseGrades(text: string): number[] {
			return text
			.split(/[,\s]+/)
			.map((token) => Number(token))
			.filter((n) => Number.isFinite(n) && n >= 0);
		}
		
		export function Students() {
			const [rows, setRows] = useState<Applicant[] | null>(null);
			const [editingId, setEditingId] = useState<number | null>(null);
			const [form, setForm] = useState(emptyForm);
			const [error, setError] = useState<string | null>(null);
			const [qrFor, setQrFor] = useState<number | null>(null);
			const [qrSvg, setQrSvg] = useState<string | null>(null);
			
			useEffect(() => {
				invoke<Applicant[]>("list_applicants").then(setRows);
			}, []);
			
			function startEdit(a: Applicant) {
				setEditingId(a.id);
				setForm({
					last_name: a.last_name,
					first_name: a.first_name,
					patronymic: a.patronymic,
					address: a.address,
					phone: a.phone,
					gradesText: a.grades.join(", "),
				});
				setError(null);
			}
			
			function cancelEdit() {
				setEditingId(null);
				setForm(emptyForm);
				setError(null);
			}
			
			async function showQr(id: number) {
				if (qrFor === id) {
					setQrFor(null);
					setQrSvg(null);
					return;
				}
				
				try {
					const svg = await invoke<string>("generate_qr", { id });
					setQrFor(id);
					setQrSvg(svg);
				} catch (err) {
					setError(String(err));
				}
			}
			
			async function remove(id: number) {
				if (!window.confirm(`Видалити абітурієнта з ID ${id}?`)) {
					return;
				}
				
				try {
					await invoke("delete_applicant", { id });
					setRows((prev) => (prev ? prev.filter((a) => a.id !== id) : prev));
					if (editingId === id) {
						cancelEdit();
					}
				} catch (err) {
					setError(String(err));
				}
			}
			
			async function submit(e: FormEvent) {
				e.preventDefault();
				const input: ApplicantInput = {
					last_name: form.last_name,
					first_name: form.first_name,
					patronymic: form.patronymic,
					address: form.address,
					phone: form.phone,
					grades: parseGrades(form.gradesText),
				};
				
				try {
					if (editingId === null) {
						const created = await invoke<Applicant>("add_applicant", { input });
						setRows((prev) => (prev ? [...prev, created] : [created]));
					} else {
						const updated = await invoke<Applicant>("update_applicant", {
							id: editingId,
							input,
						});
						setRows((prev) =>
						prev ? prev.map((a) => (a.id === editingId ? updated : a)) : prev,
						);
					}
					cancelEdit();
				} catch (err) {
					setError(String(err));
				}
			}
			
			return (
			<main className="container">
			<Link className="back" to="/">
			← На головну
			</Link>
			<h1>Абітурієнти</h1>
			
			<form className="form" onSubmit={submit}>
			<h2>{editingId === null ? "Додати абітурієнта" : `Редагувати ID ${editingId}`}</h2>
			{error && <p className="error">{error}</p>}
			<label>
			Прізвище
			<input
			required
			value={form.last_name}
			onChange={(e) => setForm({ ...form, last_name: e.currentTarget.value })}
			/>
			</label>
			<label>
			Ім'я
			<input
			required
			value={form.first_name}
			onChange={(e) => setForm({ ...form, first_name: e.currentTarget.value })}
			/>
			</label>
			<label>
			По-батькові
			<input
			value={form.patronymic}
			onChange={(e) => setForm({ ...form, patronymic: e.currentTarget.value })}
			/>
			</label>
			<label>
			Адреса
			<input
			value={form.address}
			onChange={(e) => setForm({ ...form, address: e.currentTarget.value })}
			/>
			</label>
			<label>
			Телефон
			<input
			value={form.phone}
			onChange={(e) => setForm({ ...form, phone: e.currentTarget.value })}
			/>
			</label>
			<label>
			Оцінки (через кому)
			<input
			placeholder="10, 11, 9, 12"
			value={form.gradesText}
			onChange={(e) => setForm({ ...form, gradesText: e.currentTarget.value })}
			/>
			</label>
			<div className="form-actions">
			<button type="submit">{editingId === null ? "Додати" : "Зберегти"}</button>
			{editingId !== null && (
				<button type="button" onClick={cancelEdit}>
				Скасувати
				</button>
				)}
			</div>
			</form>
			
			<section className="panel">
			<h2>
			Список <span className="count">({rows?.length ?? 0})</span>
			</h2>
			{rows && rows.length > 0 ? (
				<table>
				<thead>
				<tr>
				<th>ID</th>
				<th>ПІБ</th>
				<th>Адреса</th>
				<th>Телефон</th>
				<th>Оцінки</th>
				<th>Сер. бал</th>
				<th></th>
				</tr>
				</thead>
				<tbody>
				{rows.map((a) => (
					<tr key={a.id}>
					<td>{a.id}</td>
					<td>
					{a.last_name} {a.first_name} {a.patronymic}
					</td>
					<td>{a.address}</td>
					<td>{a.phone}</td>
					<td>{a.grades.join(", ")}</td>
					<td>{a.average}</td>
					<td>
					<button type="button" onClick={() => startEdit(a)}>
					Редагувати
					</button>
					<button type="button" onClick={() => remove(a.id)}>
					Видалити
					</button>
					<button type="button" onClick={() => showQr(a.id)}>
					{qrFor === a.id ? "Сховати QR" : "QR-код"}
					</button>
					</td>
					</tr>
					))}
				</tbody>
				</table>
				) : (
				<p className="empty">— немає —</p>
				)}
			</section>
			
			{qrFor !== null && qrSvg && (
				<section className="panel">
				<h2>QR-код абітурієнта ID {qrFor}</h2>
				<div className="qr-code" dangerouslySetInnerHTML={{ __html: qrSvg }} />
				</section>
				)}
			</main>
			);
		}
	\end{minted}
\end{document}